% !PRINTED 377
\weberpage{377}
%<<CONT
pass through the intersection of $x_1$ and $x_4$, and consequently $x_4$ is linearly composed of $x_1$ and $\lambda_2 x_3 + \dfrac{x_2}{\lambda_3}$. Now, however, we can express $x_4$ linearly in terms of $x_1, x_2, x_3$, say in the form

\weq{10}{x_4 = a_1 x_1 + a_2 x_2 + a_3 x_3,}

and in it $a_1, a_2, a_3$ are to be regarded as known, and none of these constants can vanish. Hence $\lambda_2\lambda_3 = a_3 : a_2$, and if we introduce a new constant $h_1$, then

\weq{11}{\lambda_2 = h_1 a_3,\quad \frac{1}{\lambda_3} = h_1 a_2,}

and equation (9) yields the identity:

\[
h_1 x_4\left(\frac{\xi_{42}}{\lambda_2} + \frac{\xi_{43}}{\lambda_3} - h_1\xi_1\right) = x_1\left(-h_1(2+a_1h_1)\xi_1 + \frac{\xi_2}{a_2} + \frac{\xi_3}{a_3}\right).
\]

From this one concludes further, if $k_1$ is a new constant,

\[
k_1 x_4 = -h_1(2+a_1h_1)\xi_1 + \frac{\xi_2}{a_2} + \frac{\xi_3}{a_3},
\]

\weq{12}{\frac{\xi_{42}}{\lambda_2} + \frac{\xi_{43}}{\lambda_3} = h_1\xi_1 + \frac{k_1}{h_1} x_1.}

If one interchanges the figures 1, 2, 3 cyclically in this, then three such systems result from (8); first:

\weq{13}{\begin{aligned}
k_1 x_4 &= -h_1(2+a_1h_1)\xi_1 + \frac{\xi_2}{a_2} + \frac{\xi_3}{a_3}\\
k_2 x_4 &= \frac{\xi_1}{a_1} - h_2(2+a_2h_2)\xi_2 + \frac{\xi_3}{a_3}\\
k_3 x_4 &= \frac{\xi_1}{a_1} + \frac{\xi_2}{a_2} - h_3(2+a_3h_3)\xi_3.
\end{aligned}}

Since these three expressions for $x_4$ must be identical with one another, it follows that:

\[
-\frac{h_1(2+a_1h_1)}{k_1} = \frac{1}{k_2 a_1} = \frac{1}{k_3 a_1},
\]

hence $k_2 = k_3$ and likewise $k_2 = k_1$. Thus $k_1, k_2, k_3$ are equal to one another and we write $k$ for them. Then it follows further that

\[
a_1 h_1(2+a_1h_1) + 1 = (a_1h_1+1)^2 = 0,
\]

hence $h_1 a_1 = -1$, and likewise

\[
h_1 = \frac{-1}{a_1},\quad h_2 = \frac{-1}{a_2},\quad h_3 = \frac{-1}{a_3}.
\]
